% !TEX root = ../acl_latex.tex
% Data construction pipeline figure. \input from overleaf_methods.tex inside
% a figure environment. Requires \usetikzlibrary{arrows.meta,positioning}
% in the preamble of the compiling main file. Single-column width (~8.4cm);
% laid out as a vertical chain to stay robust without a local compiler.
\centering
\begin{tikzpicture}[
  node distance=2.2mm,
  every node/.style={font=\scriptsize},
  box/.style={rectangle, draw, rounded corners, align=center,
              minimum height=6.5mm, text width=6.3cm, inner sep=2.5pt},
  arr/.style={-{Latex[length=1.8mm]}, thick}
]
\node[box] (s1) {Candidate questions: HealthCareMagic/iCliniq (PCOS, fertility) + menstrual corpus; keyword-filtered, artifact-cleaned};
\node[box, below=of s1] (s2) {Category anchored to NHS/NICHD guidance; CDC specifically for fertility where coverage was thin};
\node[box, below=of s2] (s3) {Rewritten into 6 matched registers (meaning-preservation checked)};
\node[box, below=of s3] (s4) {Risk \& clarification targets: keyword heuristic over answer text, not grounding; clarification = ambiguous-register row only};
\node[box, below=of s4] (s5) {Split: evaluation set vs.\ leakage-excluded adaptation pool};
\node[box, below=of s5] (s6) {Adaptation pool translated to French \& MSA (LLM), then QA + native-speaker confirmation};
\draw[arr] (s1) -- (s2);
\draw[arr] (s2) -- (s3);
\draw[arr] (s3) -- (s4);
\draw[arr] (s4) -- (s5);
\draw[arr] (s5) -- (s6);
\end{tikzpicture}
